\section{Data}
\label{sec:data}
% PLAN (about 1 page and one table). Physics and rendering are in methods.tex; this section lists what data are used for what, under which licence,
% and which sets are locked.
%   3.1 Training sources: LIDC-IDRI (v1 renders; lesion library for the twin pastes; synthetic test CTs) and the public Qure ct2xr-projections
%       (v2 source; exactly additive components, see Section method_v2). Licence consequence: weights trained on the projections inherit CC BY-NC-SA 4.0.
%   3.2 Real evaluation data: SZCH-X-Rays (paired DES), VinDr-RibCXR (rib masks, licensed, local only), JSRT, Node21, TBX11K, VinBigData, NIH ChestX-ray14.
%   3.3 Data discipline: development versus locked sets, split rules, what was seen before the freeze (DEVIATIONS.md 2026-10-27, SAP_v2.md).
%   Table~data: dataset, role, licence or access. Licences are copied only from pages that were read; everything else stays a \todo.
% Counts below are copied from SAP_v2.md, methods.tex and DEVIATIONS.md; regenerate them from results/ when the tables are final.
Eleven data sets serve three roles in this work: sources for training the suppressors and the CT teacher, real radiographs for evaluating
them, and an external paediatric set for testing transfer to an unseen population. Table~\ref{tab:data} lists every data set with its role
and its licence or access condition, Fig.~\ref{fig:data_splits} shows the size of every split, and Section~\ref{sec:preproc} follows each
image source from its raw form to the study that uses it.

\subsection{Training sources}
\paragraph{LIDC-IDRI} The LIDC-IDRI collection \citep{armato2011lidc,armato2015lidctcia,clark2013tcia}, with some series read from the LUNA16 copy
\citep{setio2017luna16}, supplies the v1 renders, the lesion library used to paste lesions into training images, and the synthetic test CTs. The
selection, the usable-CT rules and the patient-level split are given in Section~\ref{sec:methods}, together with the detector physics (attenuation
tables from the NIST X-ray mass attenuation coefficient database \citep{hubbell2004xraymass}, tube spectra from SpekPy \citep{poludniowski2021spekpy}).

\paragraph{Qure.ai CT projections} The v2 models are trained on the public Qure.ai \texttt{ct2xr-projections} data \citep{angaitkar2026anatomy}, which
hold, for each CT volume, a composite image and translated bone, soft-tissue and lung projections. The composite is an exact linear combination of
the components up to rounding, so the bone image of every training sample is known exactly (Section~\ref{sec:method_v2}).
The release holds \num{21887} projections, one per CT volume; the volume names are those of CT-RATE \citep{hamamci2026ctrate}. We split them by
patient (\num{14374} patients) into \num{21051} training, 406 validation and 430 test projections (\path{tools/qp_prep.py}).

\subsection{Real radiographs}
\paragraph{SZCH-X-Rays} Paired radiographs and dual-energy soft-tissue images from one partner hospital \citep{sun2026bsldm,sun2026debonedit}. We use the
75-pair test subset; its 75 images are split into a development part (24 images) and a held-out part (51 images) by a hash of the file name
(\path{tools/szch_split.py}). The 75 images are the published test split of the 741-pair version of the data set (592/74/75), acquired with dual-exposure DES on a GE
Discovery XR656 \citep{sun2026debonedit}; an earlier version of 818 pairs was used for BS-LDM \citep{sun2026bsldm}, and the papers do not
explain the difference.

\paragraph{VinDr-RibCXR} 245 radiographs with per-rib masks \citep{nguyen2021ribcxr}. The data are licensed and are processed only on the licence
holder's own computer and private cloud account; only metrics leave them.

\paragraph{Nodule and detection sets} JSRT \citep{shiraishi2000jsrt} (154 nodules), Node21 \citep{sogancioglu2024node21} (1\,476 nodule boxes,
of which two narrower than 4\,px at 512\,px have no defined contrast, so the contrast analysis uses 1\,474),
TBX11K \citep{liu2020tbx11k} and VinBigData (the Kaggle competition data; related to VinDr-CXR \citep{nguyen2022vindrcxr}) are used for real-lesion contrast
and for downstream detection; NIH ChestX-ray14 \citep{wang2017chestxray} is used for downstream classification with its official lists.
NIH uses the official test list (\num{25596} images) and a patient-level split of the remaining images into \num{77890} training and \num{8634}
validation images. The detection sets are split by a fixed hash of the image identifier: Node21 \num{3440}/\num{512}/\num{930}, TBX11K
\num{6200}/\num{688}/\num{2088} (its official validation set serves as test set) and VinBigData \num{10433}/\num{1469}/\num{2943} images for
training, validation and test. VinBigData boxes are kept where at least two radiologists marked the same finding; images with only single-reader
findings are left out. No image of any set was unreadable. VinBigData is the Kaggle competition built on the VinDr-CXR images \citep{vinbigdata2020kaggle}; its rules allow non-commercial research use.
The classifier study (Section~\ref{sec:classifier}) uses the same NIH and VinBigData splits; on VinBigData an image is positive for one of the
14 findings when at least two of its readers marked it. The 984 expert boxes on 880 NIH test images (eight classes) are used only to score
attribution maps.

\paragraph{BDCXR-3257 (external validation)} 3,257 paediatric chest radiographs (JPEG) from the Radiology and Imaging Department of
Dr.~M~R~Khan Shishu Hospital \& Institute of Child Health, Dhaka, collected in routine care and anonymised, with one label each: normal (880),
bacterial pneumonia (1,477) or viral pneumonia (900) \citep{durjoy2026bdcxr}, a task studied before on other paediatric sets \citep{kermany2018identifying}. Most children are under five years old
(1,132 images under one year). The data set is public on figshare under CC BY 4.0, but we obtained it under the ethical approval REC-FSIT-2025/No:~11928 (Daffodil
International University) and a data use agreement between the hospital and the authors' research group, whose terms we follow: no
radiograph of this set is shown, and no image or model derived from it is redistributed by us. There are no
patient identifiers, so near-duplicate images are grouped by their correlation at $64$~px ($r\ge0.93$; 48 pairs, the rest single images) and
a group never spans two folds.

\subsection{Development and locked sets}
SAP.md fixed the locked sets before any locked-set evaluation. After the v1 ablation the SZCH pairs were looked at (DEVIATIONS.md, 2026-10-27), so SZCH
became a development set for decisions about the v2 recipe, and SAP\_v2.md was written and frozen before any held-out or locked evaluation of a v2
model. Decisions used the 24 development images, the synthetic validation data and the real-image audit on all 75 SZCH images; the audit on SZCH is
therefore development evidence only, and the confirmatory audit runs on VinDr-RibCXR (not opened by any model run before the plan was written) and
JSRT. Paired fidelity claims use the 51 held-out SZCH pairs. Before the freeze, screening tools printed held-out PSNR and SSIM next to the
development numbers for v1-family and screening runs; no rule, table or decision used them, and the tool was changed to print development
numbers only.

\subsection{Data preprocessing}
\label{sec:preproc}
Every image passes through a fixed chain before any model sees it (Fig.~\ref{fig:data_pipeline}, Table~\ref{tab:preproc}):
\begin{enumerate}
\item \emph{Decomposition of CT into exact components.} LIDC-IDRI CTs are split into bone, non-lung soft tissue and lung by Hounsfield
 rules, projected along the anteroposterior axis and rendered with the polychromatic physics of Eq.~\eqref{eq:physics}; the Qure.ai
 projections are converted into an exactly additive radiograph and bone image (Eq.~\eqref{eq:qure}).
\item \emph{Augmentation that keeps the decomposition exact.} Bone re-weighting, one shared affine transform, lesion insertion and a
 monotone tone curve are applied so that $B'=X'-S'$ holds for every training sample (Eq.~\eqref{eq:compose}).
\item \emph{Real radiographs.} Each radiograph is resized to 1024\,px by area averaging and min--max scaled; the suppressors read it at
 512\,px, the classifiers at 256\,px and the detector at 512\,px.
\item \emph{Lung masks and audit sites.} Lung fields come from a published segmentation network, eroded by 49\,px, and six lesion sites per
 image are drawn once with a fixed seed, so that every model is audited at the same places (Section~\ref{sec:audit}).
\item \emph{CT volumes for Study~III.} Each CT-RATE volume is reduced to 24 axial slices over the lung extent, shown in three Hounsfield
 windows and cropped to the body at 256\,px (Section~\ref{sec:xm_data}).
\item \emph{Integrity screening and splits.} Unusable renders are removed by fixed rules, near-duplicate paediatric images are grouped, and
 every split is made by patient, image hash or duplicate group before any model sees the corresponding test data.
\end{enumerate}

\begin{figure}[htbp]
\centering
\safefig{figures/fig_data_pipeline.pdf}
\caption{Data preprocessing. Each row follows one image source from its raw form (left) through preprocessing and splitting (train /
validation / test; dev: development pairs, test: held-out pairs) to the studies that use it (right; I, II, III). Thumbnails are real images
where the licence allows: a LIDC-IDRI CT slice \citep{armato2015lidctcia} (CC~BY~3.0), a SZCH-X-Rays radiograph with its dual-energy
soft-tissue image \citep{sun2026bsldm} (CC~BY-NC~4.0) and an NIH ChestX-ray14 radiograph \citep{wang2017chestxray}. $\dagger$: no image
derived from CT-RATE is shown, so the Qure.ai row shows the composite radiograph $X$ and the exact bone image $B$ of a rendered LIDC-IDRI
CT, and the CT-RATE row its lung, mediastinum and bone windows. $\ddagger$: dummy radiographs, variants (contrast, noise, mirroring,
framing) of the same rendered LIDC-IDRI CT, stand in for data sets whose images are not shown (data use agreements or restrictive
licences: VinDr-RibCXR, JSRT, VinBigData, BDCXR-3257, Node21; TBX11K is not shown either); they are not images of these data sets.}
\label{fig:data_pipeline}
\end{figure}

\begin{figure}[htbp]
\centering
\safefig[0.6\linewidth]{figures/fig_data_splits.pdf}
\caption{Size of every data split on a logarithmic scale. SZCH-X-Rays has only development and held-out pairs, VinDr-RibCXR and JSRT are
test sets only, and BDCXR-3257 is used in five-fold cross-validation over near-duplicate groups.}
\label{fig:data_splits}
\end{figure}

\begin{table}[htbp]
\centering\small
\renewcommand{\arraystretch}{1.12}
\caption{Preprocessing of each image source. All splits are by patient, by a fixed hash of the image identifier or, for BDCXR-3257, by
near-duplicate group, and were fixed before any model saw the corresponding test data.}
\label{tab:preproc}
\begin{tabular}{@{}>{\raggedright\arraybackslash}p{0.15\linewidth}>{\raggedright\arraybackslash}p{0.16\linewidth}%
>{\raggedright\arraybackslash}p{0.38\linewidth}>{\raggedright\arraybackslash}p{0.21\linewidth}@{}}
\toprule
Source & Used by & Preprocessing & Integrity checks \\
\midrule
LIDC-IDRI CT & v1 suppressors, lesion library & Bone, non-lung soft tissue and lung masks; $\rho=\mathrm{clip}(1+\mathrm{HU}/1000,0,3)$; AP
 projection of each component; body crop, $512\times512$; physics of Eq.~\eqref{eq:physics} at 80--130~kVp (training) or 100~kVp
 (evaluation); composer of Eq.~\eqref{eq:compose} & 18 of 1\,018 series not renderable; usability rules leave 844 CTs; DICOM and LUNA16
 copies voxel-identical on 12 series\\
Qure.ai projections & v2 suppressors; Study~III $P_0$/$P_1$ & 8-bit, $512\times512$; $X=T/255$, exact $B$ (Eq.~\eqref{eq:qure}); bone
 weight log-uniform in $[0.6,1.7]$, shared affine, insertion, gamma tone curve, blur on $X'$ and $B'$, noise and scatter on $X'$ & Mean
 absolute residual of the additive fit 0.36 grey levels (largest 3.2 of 255)\\
Real radiographs (suppressor input) & Studies I and II & Suppressor run at 512\,px; lung masks from the TorchXRayVision PSPNet at
 1024\,px, eroded by 49\,px for audit sites & Masks exported once and shared by all models\\
Classifier inputs & Study~II, downstream H4$'$b, E1 & $256\times256$\,px; per-image standardisation and a learned affine map per stream;
 rotation $\le7^\circ$, scale 0.9--1.1, translation $\le3$\,\%, brightness and contrast, identical across channels & VinBigData labels
 need two readers\\
Detector inputs & Downstream H4$'$b & $512$\,px; boxes of findings marked by two or more readers & Node21: two boxes narrower than 4\,px
 excluded from contrast\\
CT-RATE volumes & Study~III teacher & 24 axial slices over the lung extent; lung, mediastinum and bone windows as channels; body crop at
 256\,px & Volumes not showing the chest or not head-first supine excluded with their projections\\
X-ray students & Study~III & $256$\,px (\texttt{swinv2\_tiny\_window8\_256}) & Same NIH and VinBigData splits as Study~II\\
BDCXR-3257 & External validation & JPEG; five class-stratified folds over near-duplicate groups (fold seed 0) & Correlation at 64\,px,
 $r\ge0.93$: 48 pairs grouped\\
\bottomrule
\end{tabular}
\end{table}

\begin{table}[htbp]
\centering\small
\renewcommand{\arraystretch}{1.15}
\caption{Data sets, their role in this study and the licence or access condition, as stated on the providers' pages (read on
3 November 2026).}
\label{tab:data}
\begin{tabular}{@{}p{0.24\linewidth}p{0.35\linewidth}p{0.33\linewidth}@{}}
\toprule
Data set & Role here & Licence or access \\
\midrule
LIDC-IDRI (TCIA) & v1 renders, lesion library, synthetic test CTs & CC BY 3.0; attribution with the data DOI \\
Qure.ai ct2xr-projections & v2 training source; Study III projection models & CC BY-NC-SA 4.0 \\
CT-RATE & source of the projections; Study III CT teacher (\num{2156} volumes) & CC BY-NC-SA 4.0, gated \\
SZCH-X-Rays & paired DES evaluation (24 development, 51 held-out) & CC BY-NC 4.0 \\
VinDr-RibCXR & locked audit and rib benchmark (245 images) & licensed; local use, metrics only \\
JSRT & audit and nodule contrast & free after registration; not public domain, no redistribution or commercial use \\
Node21 & nodule contrast, detection & CC BY-NC-ND 4.0 \\
TBX11K & detection & CC BY 4.0 (data) \\
VinBigData / VinDr-CXR & detection; classification; Study III zero-shot transfer & competition rules: non-commercial research use \\
NIH ChestX-ray14 & classification (Studies II and III) & unrestricted use with attribution (NIH Clinical Center) \\
BDCXR-3257 & external validation (pneumonia classification, audit) & CC BY 4.0 (figshare); used under the providers' data use agreement \\
\bottomrule
\end{tabular}
\end{table}
